\section{Additional $\nu_e$ $CC (1\pi^{+})$  sample at Super-K}
\label{sec:nuecconepi_sample}
Another sample in addition to the four single ring samples used in the previous VALOR joint \(\nu\)/\(\bar{\nu}\) fit analysis \cite{OurRun1617b17cJointAnalysis3Flavour} is included in the present analysis, this new sample is termed \nueCConepi{}. This is the first Super-K sample included in the analysis which includes a non CCQE-like signal topology. For this new sample the signal selection cuts allow events with one decay electron in the \(\nue\) appearance sample, whereas the single \(e\)-like ring samples in previous analyses only allowed events with no decay electrons. This new cut introduces a large \(\numu\) background, since the resulting decay electron ring may have been the result of a \(\nu_{\mu}\) CCQE interaction and subsequent \(\mu\) decay. The neutrino energy reconstruction is performed by modification of the CCQE prescription and assumes that the single pion in the final state is produced by an outgoing \(\Delta^{++}\) baryon, without considering nuclear effects. Further sample selection details, and Super-K error evaluations, can be found in \cite{InelSglPionNueDecayE}.

%%% Local Variables:
%%% mode: latex
%%% TeX-master: "jointfit_5sample2016"
%%% End:
